% Packages, colours and the small commands of the WOP manual.

\usepackage[T1]{fontenc}
\usepackage[utf8]{inputenc}
\usepackage[british]{babel}
\usepackage{lmodern}
\usepackage{microtype}
\setlength{\emergencystretch}{3em} % a line may be a little loose instead of too wide (no overfull box)
\usepackage{graphicx}
\usepackage[export]{adjustbox}
\usepackage{xcolor}
\usepackage{booktabs}
\usepackage{tabularx}
\newcolumntype{Y}{>{\raggedright\arraybackslash}X}  % a text column that is not justified
\usepackage{enumitem}
\usepackage{caption}
\usepackage[most]{tcolorbox}
\usepackage{scrlayer-scrpage}

% the version and the date come from tools/make_manual.py (build/version.tex)
\InputIfFileExists{build/version.tex}{}{\newcommand{\wopversion}{unknown}\newcommand{\wopdate}{\today}}

\usepackage[hidelinks, pdftitle={WOP user manual}, pdfsubject={Workshop Optimised Programming for CNC lathes}]{hyperref}

\definecolor{wopblue}{RGB}{40,90,170}
\setkomafont{disposition}{\color{wopblue}\sffamily\bfseries}
\setkomafont{pageheadfoot}{\sffamily\small}
\setkomafont{subject}{\color{wopblue}\sffamily\Large\bfseries}
\setkomafont{subtitle}{\color{wopblue}\sffamily\Large}
\setkomafont{title}{\color{wopblue}\sffamily\bfseries\fontsize{48}{52}\selectfont}
\setkomafont{titlehead}{\sffamily\small\color{black!60}}
\renewcommand*{\chapterheadstartvskip}{\vspace*{0.5\baselineskip}}  % a chapter starts near the top of its page
\renewcommand*{\chapterheadendvskip}{\vspace{1.2\baselineskip}}
\raggedbottom
\setcounter{tocdepth}{2}      % chapters, sections and subsections in the contents
\DeclareTOCStyleEntry[numwidth=3em]{tocline}{section}              % room for numbers like 12.10.
\DeclareTOCStyleEntry[indent=4.5em,numwidth=3.8em]{tocline}{subsection}
\setcounter{secnumdepth}{3}
\captionsetup{font=small,labelfont={sf,bf}}
\setlist{itemsep=0.2em, topsep=0.4em}

\pagestyle{scrheadings}
\clearpairofpagestyles
\ihead{\textsf{WOP user manual}}
\ohead{\headmark}
\cfoot{\pagemark}
\cfoot*{\pagemark}
\automark[section]{chapter}

% --- how things on the screen are written ---------------------------------------------------------------------------------
\newcommand{\menu}[1]{\textsf{\textbf{#1}}}                     % a menu or a menu entry
\newcommand{\menupath}[2]{\menu{#1}\,$\rightarrow$\,\menu{#2}}  % File -> Open...
\newcommand{\button}[1]{\textsf{[#1]}}                          % a button
\newcommand{\field}[1]{\textsf{#1}}                             % the name of a field or a check box
\newcommand{\key}[1]{\fbox{\small\textsf{\strut#1}}}            % a key on the keyboard
\newcommand{\file}[1]{\texttt{\detokenize{#1}}}                 % a file name or a text to type

% --- boxes ---------------------------------------------------------------------------------------------------------------
\tcbset{boxstyle/.style={breakable, boxrule=0.4pt, left=2mm, right=2mm, top=1mm, bottom=1mm, fonttitle=\sffamily\bfseries\small, before skip=0.8em, after skip=0.8em}}
\newtcolorbox{notebox}{boxstyle, colback=blue!4, colframe=blue!50!black, title=Note}
\newtcolorbox{tipbox}{boxstyle, colback=green!5, colframe=green!40!black, title=Tip}
\newtcolorbox{warningbox}{boxstyle, colback=red!5, colframe=red!60!black, title=Warning}
\newtcolorbox{draftbox}{boxstyle, colback=yellow!12, colframe=orange!70!black, title=Draft: not written yet}

% a picture of the program: \shot[width as part of the text width]{file in images/}{caption}
\newcommand{\shot}[3][0.9]{%
  \begin{figure}[htbp]\centering
  \fcolorbox{black!30}{white}{\includegraphics[max width=#1\textwidth, max height=0.62\textheight]{images/#2}}
  \caption{#3}\label{fig:#2}
  \end{figure}}
